\section{Probing Analysis: Decoding Implicit Tool Necessity}
\label{sec:probing}

\begin{table}[b!]
\vspace{-10pt}
\caption{AUROC and accuracy for predicting tool necessity from pre-generation hidden states. All probes achieve high AUROC, confirming that tool necessity is consistently encoded in models.}
\label{tab:probe_results}
\small
\centering
\begin{tabular}{lcccccc}
\toprule
& & & \multicolumn{3}{c}{AUROC by difficulty} \\
\cmidrule(lr){4-6}
Model & AUROC & Acc & Easy & Med & Hard \\
\midrule
Qwen3-1.7B            & 0.894 & 0.847 & 0.864 & 0.831 & 0.904 \\
Qwen3-4B-Inst.        & 0.948 & 0.877 & 0.933 & 0.906 & 0.948 \\
Llama-3.1-8B-Inst.    & 0.927 & 0.849 & 0.892 & 0.867 & 0.884 \\
Qwen3-14B             & 0.957 & 0.892 & 0.955 & 0.907 & 0.941 \\
Qwen3-32B             & 0.952 & 0.885 & 0.951 & 0.903 & 0.939 \\
Llama-3.3-70B-Inst.   & 0.936 & 0.872 & 0.906 & 0.849 & 0.956 \\
\bottomrule
\end{tabular}
\end{table}

We now investigate whether models encode tool-necessity information internally, even when they fail to act on it during generation. Interestingly, we find that \textit{tool necessity is already encoded in hidden states} as shown in the experiments. 


\paragraph{Setup.}
For each task, we first collect a binary label: we force the model to answer without tool access, and label tasks where it succeeds as tool-unnecessary ($y{=}0$) and tasks where it fails as tool-necessary ($y{=}1$). We then extract hidden states by running a single forward pass and taking the hidden state at the last token position across all layers. Finally, we concatenate all-layer features and train an L2-regularized logistic regression to predict tool necessity. The probe trains on 900 training examples and is evaluated on 2,250 held-out test tasks. Training takes seconds on CPU.
\vspace{-5pt}
\paragraph{Tool necessity is linearly decodable.}
The probe achieves AUROC 0.89--0.96 across all six tested models (Table~\ref{tab:probe_results}), confirming that tool necessity is consistently encoded in pre-generation hidden states regardless of model family or size. Even the smallest model carries a strong signal, while larger models reach 0.95+. The per-difficulty breakdown shows strong performance across all levels, with medium tasks being the most challenging, consistent with medium being the decision boundary where the model is uncertain.
\vspace{-5pt}
\paragraph{The signal exists even when generation fails.}
The most striking evidence comes from the Llama models. As discussed in Section~\ref{sec:behavioral}, Reason-then-Act completely breaks tool calling on Llama-3.1-8B (79.5\% $\to$ 31.2\%) and Llama-3.3-70B (83.1\% $\to$ 47.9\%). Yet the linear probe still achieves AUROC above 0.9; the information about tool necessity is clearly present in the representation, even though the model is entirely unable to express it during generation. 

Based on these findings, our next step is to use this signal to directly steer the model's tool-call behavior. In the next section, we show how injecting a short steering sentence into the model's output can translate this hidden knowledge into better tool-call decisions.
